%!TEX TS-program = pdflatex
%!TEX encoding = UTF-8 Unicode
%!TEX spellcheck = English (Aspell)

\documentclass[11pt]{article}
\usepackage{jeffe,handout,graphicx}

\def\Cdot{\mathbin{\text{\normalfont \textbullet}}}
\def\Sym#1{\textbf{\texttt{\color{BrickRed}#1}}}

\newcommand{\X}[1]{\textcolor{Red}{\mathbf{\underline{#1}}}}

\hidesolutions

% =========================================================
\begin{document}

\headers{CS/ECE 374 A}{Lab 12a — November 8}{Fall 2023}

\begin{enumerate}
\parindent 1.5em \itemsep 4ex

%----------------------------------------------------------------------
\item
Suppose you are given a magic black box that somehow answers the following decision problem \emph{in polynomial time}:
\begin{itemize}\cramped
\item \mathsc{Input:} A directed graph $G$ and a positive integer $L$.  (The edges of $G$ are not weighted, and $G$ is not necessarily a dag.)
\item \mathsc{Output:} \mathsc{True} if $G$ contains a (simple) path of length $L$, and \mathsc{False} otherwise.\footnote{You already know how to solve this problem in polynomial time \emph{when the input graph $G$ is a dag}, but this magic box works for \emph{every} input graph.}
\end{itemize}

\begin{enumerate}\itemsep 2ex
\item
Using this black box as a subroutine, describe algorithms that solves the following optimization problem \emph{in polynomial time}:
\begin{itemize}
\item \mathsc{Input:} A directed graph $G$.
\item \mathsc{Output:} The length of the longest path in $G$.
\end{itemize}

\item
Using this black box as a subroutine, describe algorithms that solves the following search problem \emph{in polynomial time}:
\begin{itemize}
\item \mathsc{Input:} A directed graph $G$.
\item \mathsc{Output:} The longest path in $G$
\end{itemize}
\end{enumerate}
\Hint{You can use the magic box more than once.}


%----------------------------------------------------------------------
\item
An \EMPH{independent set} in a graph $G$ is a subset $S$ of the vertices of $G$, such that no two vertices in $S$ are connected by an edge in $G$.  Suppose you are given a magic black box that somehow answers the following decision problem \emph{in polynomial time}:
\begin{itemize}\cramped
\item \mathsc{Input:} An undirected graph $G$ and an integer $k$.
\item \mathsc{Output:} \mathsc{True} if $G$ has an independent set of size $k$, and \mathsc{False} otherwise.\footnote{It is not hard to solve this problem in polynomial time via dynamic programming when the input graph $G$ is a \emph{tree}, but this magic box works for \emph{every} input graph.}
\end{itemize}

\begin{enumerate}\itemsep 2ex
\item
Using this black box as a subroutine, describe algorithms that solves the following optimization problem \emph{in polynomial time}:
\begin{itemize}
\item \mathsc{Input:} An undirected graph $G$.
\item \mathsc{Output:} The size of the largest independent set in $G$.
\end{itemize}
% \Hint{You’ve seen this problem before.}

\item
Using this black box as a subroutine, describe algorithms that solves the following search problem \emph{in polynomial time}:
\begin{itemize}
\item \mathsc{Input:} An undirected graph $G$.
\item \mathsc{Output:} An independent set in $G$ of maximum size.
\end{itemize}
\end{enumerate}
\Hint{You can use the magic box more than once.}


\end{enumerate}

%----------------------------------------------------------------------
\newpage

\noindent
\textbf{To think about later:}

\begin{enumerate}\parindent 1.5em\itemsep4ex
\setcounter{enumi}{2}

%----------------------------------------------------------------------
\item
Suppose you are given a magic black box that somehow answers the following decision problem in \emph{polynomial time}:
\begin{itemize}
\item
\mathsc{Input:} A boolean circuit $K$ with $n$ inputs and one output.
\item
\mathsc{Output:} \mathsc{True} if there are input values $x_1, x_2, \dots, x_n \in \set{\mathsc{True},\mathsc{False}}$ that make $K$ output \mathsc{True}, and \mathsc{False} otherwise.
\end{itemize}
Using this black box as a subroutine, describe an algorithm that solves the following related search problem \emph{in polynomial time}:
\begin{itemize}
\item \mathsc{Input:} A boolean circuit $K$ with $n$ inputs and one output.
\item \mathsc{Output:} Input values $x_1, x_2, \dots, x_n \in \set{\mathsc{True},\mathsc{False}}$ that make $K$ output \mathsc{True}, or \mathsc{None} if there are no such inputs.
\end{itemize}
\Hint{You can use the magic box more than once.}


\end{enumerate}



\end{document}